\pagenumbering{arabic} %necessary to start numbering pages in arabic numerals from here
\setcounter{page}{1} 

\section{Introducción}
Los problemas de optimización aparecen en prácticamente todas las áreas de la ingeniería y la ciencia, desde el ajuste de parámetros en modelos de aprendizaje automático hasta el diseño de estructuras y el control de sistemas dinámicos. Muchos de estos problemas no tienen una solución analítica sencilla: la función objetivo puede ser no lineal, no convexa, o carecer de una expresión derivable que permita aplicar métodos clásicos basados en gradiente. En estos casos, las metaheurísticas bioinspiradas ofrecen una alternativa práctica, ya que buscan soluciones aceptables sin requerir información detallada sobre la estructura del problema.
\\\\
La Optimización por Enjambre de Partículas (PSO, por sus siglas en inglés) es una de estas metaheurísticas. Fue propuesta por Kennedy y Eberhart en 1995 a partir de la observación del comportamiento colectivo de bandadas de aves y bancos de peces, donde cada individuo ajusta su trayectoria con base en su propia experiencia y en la información compartida por el grupo. Esta idea se tradujo en un algoritmo donde un conjunto de partículas explora el espacio de búsqueda, actualizando su posición y velocidad en función de su mejor resultado histórico y del mejor resultado encontrado por el enjambre.
\\\\
El comportamiento del algoritmo depende en gran medida de sus parámetros de control, en particular del peso de inercia, que regula el equilibrio entre exploración y explotación durante la búsqueda. En su formulación original, este peso permanece constante a lo largo de todas las iteraciones. Sin embargo, trabajos posteriores mostraron que variar la inercia de forma dinámica, ya sea disminuyéndola o aumentándola conforme avanza el proceso, puede mejorar la capacidad del algoritmo para converger hacia mejores soluciones.
\\\\
En esta práctica se implementa primero la versión clásica de PSO, con parámetros fijos, y después se extiende en dos pasos. Primero se hace variable únicamente el peso de inercia $w$, con una estrategia lineal decreciente y otra basada en el fitness de cada partícula. Después se hacen variables los tres coeficientes $w$, $c_1$ y $c_2$, siguiendo el esquema de coeficientes de aceleración variables en el tiempo (TVAC). El objetivo es comparar estos enfoques y analizar el efecto que tiene cada modificación sobre el comportamiento del enjambre y la calidad de las soluciones obtenidas.